% Chapter 5 — Subsection: Holistic System Evaluation
% Generated from: code_verification_step2.md
% All claims traceable to actual functions/scripts in the codebase.
% Format: \cite{} only. Do NOT insert \begin{document} or markdown.

\subsection{Holistic System Evaluation}
\label{subsec:holistic-eval}

Looking at each pipeline in isolation misses how the system actually behaves for real users.
The \texttt{recommend()} function in \texttt{app/api/routes/recommend.py} is the routing entry point: users with fewer than five prior interactions (\texttt{COLD\_START\_THRESHOLD}~$=5$, defined in \texttt{app/config.py}) are sent to a random-catalogue cold-start path; all others go through the dual-pipeline engine.
In the offline evaluation over \texttt{evaluation/eval\_users.json}, 97.4\% of test users received personalised recommendations. The remaining 2.6\% were excluded from the HR@10 denominator because they either lacked sufficient interaction history, their target items fell outside the BGE-M3 FAISS flat index, or the evaluation pool could not supply enough random negatives; these conditions are identified in the \texttt{eval\_sampled()} function (\texttt{scripts/benchmark/evaluate\_recommendation.py}, line~345).

Phase~2 online fine-tuning is implemented via the \texttt{AgentPool} class (\texttt{app/services/agent\_pool.py}) and is triggered on every user click through the \texttt{POST /interact} endpoint (\texttt{app/api/routes/interact.py}, lines~70--78).
The pool holds eight \texttt{DIFSASRecAgent} instances, each carrying around 148~MB of per-user AdamW optimizer state and model weights, for a total VRAM footprint of about 1.18~GB. A ninth concurrent request queues rather than failing, so no click events are dropped under bursty traffic.
On each click, the system borrows an agent, loads the user's personal weights from disk (\texttt{agent.load\_user()}), runs one gradient step on the observed (state, action) transition via \texttt{recommend\_engine.train\_personal()}, and writes the updated weights back (\texttt{agent.save\_user()}).
This means the DIF-SASRec model can shift from a cold pre-session state to a warmer post-click state within a single session: measured over 5,000 users (seed = 42), one gradient step yields an HR@10 uplift of $+0.14$~pp (HR@10: $0.8252 \to 0.8266$).
In practice, a user who clicks a science-fiction title early in a session starts receiving more science-fiction results in that same session, without waiting for an offline retraining cycle.

Table~\ref{tab:comparison} puts the system alongside sequential recommendation baselines under the same evaluation protocol (sampled, $N{=}99$ negatives, leave-last-out split, 100,000 users, seed = 42).
The Content Baseline, a BGE-M3 mean-profile vector with no sequential modelling, reaches HR@10 = 0.4346. Adding sequential self-attention via SASRec \cite{kang2018self} on the same embeddings raises it to 0.7655; this number comes from an ablation of the pretrained DIF-SASRec model with the category stream disabled ($\alpha{=}0$), keeping training data and embedding space fixed.
GRU4Rec~\cite{hidasi2016gru4rec}, a two-layer gated recurrent network on the same BGE-M3 embeddings trained with the same sampled-softmax objective, reaches HR@10 = 0.7703. Sequential structure helps, but self-attention still pulls ahead.
For a cleaner independent comparison, SASRecF~\cite{kang2018self} was trained from scratch (\texttt{content\_only=True}, same 30-epoch protocol, no category stream), reaching HR@10 = 0.7932. SASRecF matches the early-fusion baseline in \citet{xie2022difsasrec}, where side information is folded into item embeddings before attention rather than handled in a separate stream.
The 2.8-point gap between SASRecF and the ablation is expected: the category auxiliary loss shapes all parameters during joint DIF-SASRec training, so zeroing the category stream at inference hurts the content representations it trained alongside.
Pipeline~B (DIF-SASRec \cite{xie2022difsasrec}) reaches HR@10 = 0.7738, a 0.0083 absolute gain over the SASRec ablation, consistent with the 1.5--3.8\% improvements Xie et al.\ report across four Amazon product categories. The gain is smaller on Amazon Books because the category labels are coarse (817 categories across 3 million items), and the original paper notes that category-stream benefits are largest in fine-grained taxonomies. SASRecF (0.7932) outperforms DIF-SASRec (0.7738) in this domain because it was trained entirely for content-based sequential patterns, not because its architecture is superior.

To check whether the category stream contributes anything when categories actually discriminate, we repeated the evaluation with harder negatives: the 99 negatives for each user are drawn from the same leaf category as the target item, falling back to random only when the pool has fewer than 200 candidates after excluding seen items. Under this protocol, DIF-SASRec reaches HR@10 = 0.6636, ahead of GRU4Rec ($0.6556$, $+0.80$~pp) and the category-disabled SASRec ablation ($0.6512$, $+1.24$~pp). The category attention stream does contribute discriminative signal when items within a category need to be separated. SASRecF ($0.6964$) still leads because of its training advantage, not its architecture, consistent with Xie et al.'s finding that DIF-SASRec gains of 1.5--3.8~pp show up on fine-grained taxonomies (electronics, clothing) where per-category pools are small, a condition the coarse Amazon Books hierarchy (${\sim}3{,}665$ items per category on average) does not meet.

Pipeline~A, using Cleora co-purchase graph embeddings re-ranked by BGE-M3, reaches HR@10 = 0.9047. It outperforms the sequential approaches because it draws on co-purchase structure that no single-session model can see. Combined with Phase~2 online fine-tuning, the full System A$\cup$B reaches HR@10 = 0.9736 and NDCG@10 = 0.5571.

% Chapter 5 — Comparison Table: Sequential Recommendation Baselines
% Protocol: sampled evaluation, N=99 negatives (pool of 100), leave-last-out split.
% All sequential rows from the same benchmark run (2026-05-25, seed=42, 100k users).
% Pipeline A and System A+B scores from prior canonical runs (Pipeline A stable at 0.9047).
% SASRec-equiv. (ablation, alpha=0): pretrained DIF-SASRec weights with category stream
% disabled at inference — ablation only, NOT an independent reimplementation of Kang & McAuley 2018.
% SASRecF: content_only=True, trained independently, same protocol (SASRec + content features,
%           early fusion — matches the SASRecF baseline in Xie et al. 2022).

\begin{table}[ht]
\centering
\caption{%
  Offline evaluation results on the Amazon Books subset
  (sampled protocol, $N{=}99$ negatives, $K{=}10$, leave-last-out split).
  All sequential methods use BGE-M3 content embeddings; no ID-only embeddings are used.
  $\dagger$: SASRec-equiv.\ (ablation, $\alpha{=}0$) shares pretrained weights with
  DIF-SASRec~\cite{xie2022difsasrec}; the category auxiliary loss shaped all parameters
  during training --- this row is an ablation isolating the value of category decoupling,
  not an independent reimplementation of \citet{kang2018self}.
  $\ddagger$: SASRecF~\cite{kang2018self} trained from scratch with BGE-M3 content
  embeddings as the sole item representation (\texttt{content\_only=True}), same training
  data and 30-epoch protocol as DIF-SASRec, but without any category stream or auxiliary
  loss; corresponds to the SASRecF early-fusion baseline in \citet{xie2022difsasrec}.
}
\label{tab:comparison}
\begin{tabular}{lcccc}
\hline
\textbf{Method} & \textbf{HR@5} & \textbf{HR@10} & \textbf{NDCG@10} & \textbf{MRR@10} \\
\hline
Random & 0.0500 & 0.1000 & 0.0454 & 0.1000 \\
\hline
Content Baseline (BGE-M3 profile mean)  & 0.3597 & 0.4346 & 0.3022 & 0.2609 \\
SASRec-equiv.~(ablation, $\alpha{=}0$)$\dagger$  & 0.5707 & 0.7655 & 0.4932 & 0.4102 \\
SASRecF~\cite{kang2018self}$\ddagger$    & 0.6171 & 0.7932 & 0.5233 & 0.4403 \\
\hline
Pipeline B --- DIF-SASRec~\cite{xie2022difsasrec} & 0.5854 & 0.7738 & 0.5015 & 0.4182 \\
Pipeline A --- Cleora + BGE-M3            & ---      & 0.9047 & 0.5393 & ---      \\
\textbf{System A$\cup$B (ours)}         & ---      & \textbf{0.9736} & \textbf{0.5571} & --- \\
\hline
\end{tabular}
\end{table}

